% Generated draft: requires mathematical and source review.
\documentclass[11pt]{article}
\usepackage[margin=25mm]{geometry}
\usepackage{fontspec}
\setmainfont{Latin Modern Roman}
\usepackage{amsmath,amssymb,mathtools,unicode-math}
\setmathfont{Latin Modern Math}
\usepackage{enumitem,needspace,xurl,hyperref}
\setlength{\parindent}{0pt}
\setlength{\parskip}{5pt}
\setlength{\emergencystretch}{3em}
\newcommand{\N}{\mathbb{N}}
\newcommand{\Z}{\mathbb{Z}}
\newcommand{\Q}{\mathbb{Q}}
\newcommand{\R}{\mathbb{R}}
\newcommand{\C}{\mathbb{C}}
\newcommand{\F}{\mathbb{F}}
\newcommand{\A}{\mathbb{A}}
\newcommand{\PP}{\mathbb{P}}
\newcommand{\E}{\mathbb{E}}
\newcommand{\eps}{\varepsilon}
\newcommand{\dd}{\,\mathrm{d}}
\begin{document}
\section*{UnsolvedMath ID: 30005016}
\textbf{Source problem code:} catalogue entry
\subsection{Definitions and mathematical statement}
Let \(K\) be a field. The supplied record does not specify the field or explicitly identify the domain of the planar points. Under the interpretation that the points belong to \(K^2\), the interpolation question can be stated as follows.

For a four-element subset \(X\subseteq K^2\), let \(K^X\) denote the vector space of all functions from \(X\) to \(K\). For a four-dimensional \(K\)-linear subspace \(V\subseteq K[x,y]\), define the evaluation map
\[
\operatorname{ev}_{X,V}:V\longrightarrow K^X,
\qquad p\longmapsto \bigl(a\mapsto p(a)\bigr).
\]
Say that \(V\) interpolates arbitrary data on \(X\) when this map is surjective. Since both spaces have dimension four, this is equivalent to the evaluation map being an isomorphism.

Determine the minimum size of a family \(\mathcal F\) of four-dimensional \(K\)-linear subspaces of \(K[x,y]\) satisfying
\[
\forall X\subseteq K^2,\quad |X|=4
\ \Longrightarrow\ 
\exists V\in\mathcal F\ \forall f:X\to K\ \exists p\in V\ 
\forall a\in X,\quad p(a)=f(a).
\]
Here \(\mathcal F\) is chosen independently of \(X\), while \(V\) may depend on \(X\) and must interpolate every function on that same set. The points are distinct because \(X\) is a set of cardinality four. No degree bound on the members of \(V\) is supplied. A numerical target requires clarification of \(K\), the intended point domain, and whether the minimum is sought among finite families or as a cardinality allowing infinite families. This formulation is conditional on the stated interpretation of the point domain.

The original extract begins with the word ``corresponding'', whose antecedent is absent. It then appears to run into a separate problem headed ``Dimensions of Spline spaces'', attributed to Schenck. That material is retained here as a distinct source qualification, not as an additional hypothesis on the interpolation family. Let \(\Delta\) be a simplicial subdivision (triangulation) of a simply connected bounded domain in \(\R^2\). For positive integers \(r,m\), let
\[
S_m^r(\Delta)=\{s:s\text{ is }C^r\text{ on the domain and }s|_T
\text{ is a polynomial of total degree at most }m
\text{ for every triangle }T\in\Delta\}.
\]
This is a real vector space. The adjacent question asks for its dimension in terms of the combinatorics, local geometry, and global geometry of \(\Delta\). The extract describes this question as open even in a case printed as \(\texttt{S31}\), apparently \(S_3^1(\Delta)\); the indices require confirmation. It also describes a long history and the algebraic smoothness condition: for adjacent triangles with polynomial pieces \(f,g\), and an affine linear polynomial \(\ell\) whose zero line contains their common edge, matching derivatives through order \(r\) across the edge is expressed by
\[
\ell^{r+1}\mid(f-g),\qquad\text{equivalently}\qquad
f-g\in\langle\ell^{r+1}\rangle.
\]
These spline assertions and their open-status description are supplied source content, not independently verified claims.
\subsection{Short English statement}
Find the smallest fixed family of four-dimensional polynomial spaces such that, for every set of four distinct planar points, one space in the family interpolates all prescribed values on those points. The field and point domain remain unspecified; the displayed formulation assumes points in \(K^2\). The original extract also contains an apparently separate spline-dimension problem attributed to Schenck.
\subsection{Sources}
\begin{itemize}
\item Supplied local UnsolvedMath registry record, permanent ID \(30005016\), title ``Universal Four-Point Polynomial Interpolation Subspaces'', exact source problem code \texttt{catalogue entry}; collection: Oberwolfach Reports Open Problems. No proposer is supplied for the interpolation question. The registry gives proposed year \(2022\).
\item Supplied source citation: \emph{Mini-Workshop: Interpolation, Approximation, and Algebra (2023)}, Oberwolfach Reports, pp.~383--428, DOI \texttt{10.4171/owr/2022/7}. Source URL: \url{https://doi.org/10.4171/owr/2022/7}. The supplied title year \(2023\) differs from the DOI year component and proposed year \(2022\); this discrepancy is unresolved.
\item Supplied literature source, labelled ``OWR interpolation report'': \url{https://doi.org/10.4171/owr/2022/7}. The exact supplied literature-check date is \texttt{2026-08-22}. This is supplied dated curation, not an independently verified literature assessment. It labels the interpolation question open and reports that no general determination of the minimum was found. The background describes the review as dated literature triage derived from OWR curation and requests independent verification of time-sensitive status and cited claims.
\item The registry labels the statement \texttt{corrected\_verified} and says it was checked against the official report and preserves its hypotheses, notation, and target. These are supplied verification assertions; this draft does not independently confirm them. The original extract additionally attributes the adjacent spline-dimension problem to Schenck.
\end{itemize}
\end{document}
